%%%%%%%%%%%%%%%%%%%%%%%%%%%%%%%%%%%%%%%%%%%%%%%%%%%%%%%%%%%%%%%%%%%%%%%%%%%%%%%%
\section{Efficiency Profile of Graph Reduction Techniques (RQ2)}
\label{sec:results:rq2}
% Answers \ref{sec:intro:rq2}: shrinkage, memory, and speed across all reduction methods.

\subsection{Graph Reduction Offline Profile}
\label{sec:results:rq2:offline}
How much each method removes, and what it costs to remove it. Node and edge retention are
reported separately, because several methods preserve node counts exactly and a single
compression figure would hide that (\ref{sec:prelim:reduction:measuring}).

\begin{figure*}[htbp]
    \centering
    \includegraphics[width=\textwidth]{\resultsfigs/rq2_retention}
    \caption[Node and edge retention]{\textbf{PARTLY FABRICATED.} Offline compression per
    method, reported separately for nodes and edges. The four partitioners keep every node
    by construction, as do two of the four sparsifiers. The summarization rows are invented,
    since no summarization configuration has been trained or measured.}
    \label{fig:rq2_retention}
\end{figure*}

Retention is the axis every later result is read against, and the spread across it is what
makes a raw ranking of methods meaningless on its own: METIS keeps 93.7\% of edges while
random hashing keeps 40.2\%, so their accuracies in RQ3 are not comparable as-is.

\begin{figure*}[htbp]
    \centering
    \includegraphics[width=\textwidth]{\resultsfigs/rq2_retention_distribution}
    \caption[Retention distributions]{Per-graph retention over 10{,}000 graphs. Compression
    is a distribution rather than a ratio: for the parameter-free methods it is a property
    of each circuit, and PageRank pruning's node retention is nearly a point mass at its
    configured 0.8 while AND-gate-only's is not.}
    \label{fig:rq2_retention_distribution}
\end{figure*}

The width of these distributions is what makes the RQ4 pairing between AND-gate-only and
PageRank pruning approximate rather than exact
(\ref{sec:results:rq4:pairings}). A method whose compression varies with the
circuit cannot be matched to a fixed ratio, only calibrated to one on average.

\begin{figure}[htbp]
    \centering
    \includegraphics[width=\columnwidth]{\resultsfigs/rq2_offline_cost}
    \caption[Offline cost per method]{\textbf{PARTLY FABRICATED.} Wall-clock cost of
    computing each reduction, per graph. The measured methods span two orders of magnitude,
    from 0.004\,s for random hashing to 3.7\,s for span-weighted METIS. The summarization
    entries are invented, since no summarization configuration has been run.}
    \label{fig:rq2_offline_cost}
\end{figure}

Offline cost is paid once per graph and reused by every epoch that reads the cached
artifact, which is why it is reported as an amortisation threshold rather than as a raw
number of seconds.

\begin{figure}[htbp]
    \centering
    \includegraphics[width=\columnwidth]{\resultsfigs/rq2_amortisation}
    \caption[Amortisation threshold]{How many epochs over the corpus a cached reduction must
    serve before its offline cost is repaid by the training time it saves. This is the form
    in which~\ref{sec:method:experiment:metrics:reduction} asks offline cost to be
    reported.}
    \label{fig:rq2_amortisation}
\end{figure}

A reduction pays for itself once the epochs served exceed $c_R / s$, the offline cost
divided by the per-epoch saving. Spanning forest has no such threshold at all: its measured
time saving is $-0.1\%$, so no number of epochs repays its 3.3\,s per graph.

\begin{figure}[htbp]
    \centering
    \includegraphics[width=\columnwidth]{\resultsfigs/rq2_partition_balance}
    \caption[Partitioner trade-off]{Edge cut against partition imbalance for the four
    partitioners at the same $k$. Since all four keep every node and differ only in which
    edges they cut, this plane is their entire comparison.}
    \label{fig:rq2_partition_balance}
\end{figure}

\begin{table*}[htbp]
\centering
\caption{The four partitioners keep every node and differ only in which edges they cut, at the same $k$. That makes them the cleanest domain-informed / generic pairs in the thesis.}
\label{tab:rq2_partition_balance}
\small
\resizebox{\textwidth}{!}{%
\begin{tabular}{llrrrrrr}
\toprule
\textbf{Partitioner} & \textbf{Kind} & \textbf{Mean $k$} & \textbf{Edge cut} & \textbf{Edge cut sd} & \textbf{Imbalance (sd nodes)} & \textbf{Max part. nodes} & \textbf{Offline (s)} \\
\midrule
Random hashing & Generic & 4.662 & 0.598 & 0.162 & 37.891 & 6382 & 0.004 \\
METIS & Generic & 4.662 & 0.063 & 0.053 & 30.433 & 6370 & 3.196 \\
Level slicing & Domain & 4.662 & 0.295 & 0.236 & 0.378 & 6334 & 1.151 \\
Span-weighted METIS & Domain & 4.662 & 0.077 & 0.060 & 30.616 & 6371 & 3.701 \\
\bottomrule
\end{tabular}
}
\end{table*}


Edge cut and imbalance are the two quantities a partitioner trades against each other, and
holding $k$ fixed across all four reduces the comparison to that single plane. The pairs it
defines are the cleanest domain-informed against generic comparisons in the thesis, and RQ4
uses them as such (\ref{sec:results:rq4:pairings}).

\begin{figure}[htbp]
    \centering
    \includegraphics[width=\columnwidth]{\resultsfigs/summ_wl_depth_probe}
    \caption[Residual redundancy after strashing]{\textbf{ENTIRELY FABRICATED. NO
    SUMMARIZATION DATA EXISTS.} Node retention under colour refinement at depths
    $d = 1 \dots 4$. Every corpus graph is structurally hashed, so depth~1 should find
    almost nothing; how much survives beyond one hop bounds what the exact summarization
    track can compress and is a reportable statistic about AIGs in its own right.}
    \label{fig:summ_wl_depth_probe}
\end{figure}

The probe is the cheapest of the outstanding summarization measurements, since it needs no
training run, only an offline pass. Until it exists there is no upper bound on what exact
coarsening could remove from this corpus.
\clearpage

\subsection{Training Memory Dynamics and Compute Efficiency}
\label{sec:results:rq2:memory}
Memory is reported as peak allocated device memory per configuration, in absolute terms and
as a fraction of device capacity. Compute efficiency is reported as device utilisation
alongside memory-bandwidth activity, which together distinguish a configuration that is
compute-bound from one that is waiting on memory.
% Host-side memory is logged; report only if it constrains a configuration.

\begin{figure*}[htbp]
    \centering
    \includegraphics[width=\textwidth]{\resultsfigs/rq2_vram}
    \caption[Peak training memory]{Peak allocated device memory per configuration at three
    order statistics. The mean, p95 and maximum disagree by an order of magnitude: a graph
    larger than the batch budget forms a batch of its own, so the maximum is floored by the
    largest circuit in the corpus no matter what the reduction did to the average one.}
    \label{fig:rq2_vram}
\end{figure*}

Peak memory tracks \emph{node} retention, not edge retention. Activation memory in message
passing scales with the number of nodes times the number of layers times the hidden width,
whereas removing an edge only shortens the gather and scatter that feed those activations,
which is transient work rather than stored state. The measured savings follow that exactly.
Random edge dropout and spanning forest hold
$|\mathcal{V}(R(G))| / |\mathcal{V}(G)| = 1$ and save 1.5\% and 2.1\% of mean peak memory
while discarding 30.3\% and 41.9\% of edges respectively, and the four partitioners, which
also keep every node, save between 0.3\% and 3.1\%. The two methods that actually remove
nodes are the only ones with a double-digit saving: PageRank pruning at 0.800 node retention
saves 14.5\%, and AND-gate-only at 0.821 saves 11.3\%. Reading the maximum instead of the
mean flattens even that difference, because the largest circuit forms its own batch and
floors every configuration alike. Sparsification, in short, is not a memory intervention on
this encoder; it is a bandwidth intervention that happens to leave the memory bill intact.

\begin{figure*}[htbp]
    \centering
    \includegraphics[width=\textwidth]{\resultsfigs/rq2_cost_vs_size}
    \caption[Cost profile against graph size]{Peak memory and step time as a function of
    graph size. In absolute terms every method traces the same memory curve, because memory
    tracks the graph rather than the reduction, so the two ratio panels are where any
    difference between methods is visible at all.}
    \label{fig:rq2_cost_vs_size}
\end{figure*}

Reporting a single peak per configuration would conflate the reduction's effect with which
graphs happened to land in the peak batch, which is why the cost profile is given against
size rather than as a scalar.
\clearpage

\subsection{Throughput and Training Speedup}
\label{sec:results:rq2:throughput}
Step time and epoch time against the full-graph baseline, and inference throughput in graphs
per second over the same reference.

\begin{figure*}[htbp]
    \centering
    \includegraphics[width=\textwidth]{\resultsfigs/rq2_throughput}
    \caption[Training throughput]{Mean step time and speedup against the full-graph
    baseline. Speedup under-delivers relative to compression, because the kernels are
    latency-bound on message-passing gather and scatter
    (\ref{sec:method:experiment:hardware}), so removing work does not translate
    proportionally into wall-clock time.}
    \label{fig:rq2_throughput}
\end{figure*}

The mismatch is large enough to invert the ranking. Spanning forest removes 41.9\% of edges
and returns no time at all, while METIS removes 6.3\% and returns 24.3\%, because the
partitioners cut the batch into pieces the kernels schedule better rather than simply giving
them less to do.

\begin{figure*}[htbp]
    \centering
    \includegraphics[width=\textwidth]{\resultsfigs/rq2_paired_savings}
    \caption[Paired per-graph savings]{Per-graph peak-memory and step-time savings against
    the baseline, paired by graph id, with 95\% percentile bootstrap intervals. A red
    $p$-value marks a Wilcoxon signed-rank test that fails to separate the saving from
    noise.}
    \label{fig:rq2_paired_savings}
\end{figure*}

\begin{table*}[htbp]
\centering
\caption{Per-graph savings against the full-graph baseline, paired by graph id, with percentile bootstrap intervals and a Wilcoxon signed-rank test on the paired deltas. Savings are not normally distributed, so bare means would be misleading.}
\label{tab:rq2_paired_savings}
\small
\resizebox{\textwidth}{!}{%
\begin{tabular}{lrrrrrrrr}
\toprule
\textbf{Method} & \textbf{Paired graphs} & \textbf{VRAM saving (\%)} & \textbf{95\% CI low} & \textbf{95\% CI high} & \textbf{Time saving (\%)} & \textbf{95\% CI low } & \textbf{95\% CI high } & \textbf{Wilcoxon $p$ (time)} \\
\midrule
PageRank pruning & 1000 & 14.492 & 14.121 & 14.842 & 19.145 & 18.081 & 20.203 & 0.000 \\
AND-gate only & 1000 & 11.257 & 10.702 & 11.794 & 28.893 & 27.627 & 30.054 & 0.000 \\
Random hashing & 1000 & 3.100 & 2.984 & 3.216 & 20.996 & 19.785 & 22.157 & 0.000 \\
Spanning forest & 1000 & 2.142 & 2.080 & 2.202 & -0.101 & -1.031 & 0.826 & 0.549 \\
Random edge dropout & 1000 & 1.489 & 1.451 & 1.526 & 0.229 & -0.574 & 1.012 & 0.001 \\
Level slicing & 1000 & 1.487 & 1.383 & 1.596 & 18.291 & 17.322 & 19.300 & 0.000 \\
Span-weighted METIS & 1000 & 0.357 & 0.336 & 0.378 & 22.412 & 21.396 & 23.435 & 0.000 \\
METIS & 1000 & 0.279 & 0.262 & 0.297 & 24.284 & 23.237 & 25.303 & 0.000 \\
\bottomrule
\end{tabular}
}
\end{table*}


Savings are not normally distributed, so bare means would be misleading and the paired
bootstrap is reported instead. Spanning forest's step-time saving sits inside the noise
despite removing 42\% of edges, which is the clearest single statement of the gap between
compression and speed.
\clearpage

\subsection{Efficiency Ranking Across Methods}
\label{sec:results:rq2:ranking}
The same measurements as above, cut the other way: one row per method, ranked on what it
buys. Generic against domain-informed is a column rather than the organising axis, because
it is one attribute of a method.

\begin{table*}[htbp]
\centering
\caption{RQ2 efficiency profile. Node and edge retention are reported separately because several methods preserve node counts exactly. Savings are per-graph means paired against the full-graph baseline.}
\label{tab:rq2_efficiency}
\small
\resizebox{\textwidth}{!}{%
\begin{tabular}{lllrrrrrrr}
\toprule
\textbf{Method} & \textbf{Family} & \textbf{Kind} & \textbf{Node ret.} & \textbf{Edge ret.} & \textbf{Offline (s/graph)} & \textbf{Step time (s)} & \textbf{VRAM mean (MB)} & \textbf{VRAM saving (\%)} & \textbf{Time saving (\%)} \\
\midrule
Random hashing & Partition & Generic & 1.000 & 0.402 & 0.004 & 0.026 & 1002.217 & 3.100 & 20.996 \\
METIS & Partition & Generic & 1.000 & 0.937 & 3.196 & 0.027 & 1056.507 & 0.279 & 24.284 \\
Level slicing & Partition & Domain & 1.000 & 0.705 & 1.151 & 0.029 & 1019.120 & 1.487 & 18.291 \\
Span-weighted METIS & Partition & Domain & 1.000 & 0.923 & 3.701 & 0.027 & 1055.088 & 0.357 & 22.412 \\
Random edge dropout & Sparsification & Generic & 1.000 & 0.697 & 0.064 & 0.040 & 1040.969 & 1.489 & 0.229 \\
Spanning forest & Sparsification & Generic & 1.000 & 0.581 & 3.309 & 0.041 & 1031.674 & 2.142 & -0.101 \\
PageRank pruning & Sparsification & Generic & 0.800 & 0.627 & 1.739 & 0.029 & 850.969 & 14.492 & 19.145 \\
AND-gate only & Sparsification & Domain & 0.821 & 0.733 & 0.099 & 0.026 & 892.163 & 11.257 & 28.893 \\
{[TODO/FAKE] Random within-type} & Summarization & Generic & 0.700 & 0.700 & 0.500 & 0.030 & 800.000 & -- & -- \\
{[TODO/FAKE] ConvMatch} & Summarization & Generic & 0.500 & 0.600 & 20.000 & 0.028 & 700.000 & -- & -- \\
{[TODO/FAKE] Cone coarsening} & Summarization & Domain & 0.600 & 0.650 & 2.000 & 0.026 & 750.000 & -- & -- \\
{[TODO/FAKE] MFFC coarsening} & Summarization & Domain & 0.400 & 0.500 & 3.000 & 0.024 & 600.000 & -- & -- \\
{[TODO/FAKE] Colour refinement (exact)} & Summarization & Domain & 0.800 & 0.850 & 5.000 & 0.032 & 900.000 & -- & -- \\
\bottomrule
\end{tabular}
}
\\[4pt]\begin{minipage}{\linewidth}\footnotesize\emph{Rows marked [TODO/FAKE] are invented. No summarization configuration has been trained or measured. Needs: sbatch src/shell/precompute\_summarization.sh for the offline pass, then one training run and one src/test.py pass per method. Random within-type merging is not implemented at all.}\end{minipage}
\end{table*}


Read down the columns and no method wins on all of them. AND-gate-only returns the largest
time saving at 28.9\% and the second-largest memory saving at 11.3\% for 0.099\,s of offline
work per graph, which makes it the best-value reduction measured here on efficiency alone;
RQ3 then asks what it costs in accuracy. The summarization rows are invented and carry no
efficiency measurement at all, so the ranking covers partitioning and sparsification only.
\clearpage
